\documentclass[12pt,a4paper]{report}
\usepackage[utf8]{inputenc}
\usepackage[T1]{fontenc}
\usepackage[spanish,es-nodecimaldot]{babel}
\usepackage[a4paper,margin=2.45cm,headheight=26pt]{geometry}
\usepackage{lmodern}
\usepackage{graphicx}
\usepackage{xcolor}
\usepackage{booktabs}
\usepackage{longtable}
\usepackage{listings}
\usepackage{tikz}
\usetikzlibrary{positioning,arrows.meta,calc}
\usepackage{float}
\usepackage{caption}
\usepackage{tcolorbox}
\usepackage{fancyhdr}
\usepackage{enumitem}
\usepackage{hyperref}

\definecolor{techpurple}{HTML}{7135A5}
\definecolor{techpurple2}{HTML}{914BC6}
\definecolor{techlavender}{HTML}{F3EFF8}
\definecolor{techink}{HTML}{25212B}
\definecolor{techpink}{HTML}{D94E91}
\definecolor{techline}{HTML}{DED0EA}
\hypersetup{colorlinks=true,linkcolor=techpurple,urlcolor=techpurple2,citecolor=techpurple}
\pagestyle{fancy}
\fancyhf{}
\lhead{\textcolor{techpurple}{\textbf{TECH-FLOW}}}
\rhead{\textcolor{techpurple}{\leftmark}}
\cfoot{\textcolor{techpurple}{\thepage}}
\renewcommand{\headrulewidth}{0.5pt}
\renewcommand{\headrule}{\hbox to\headwidth{\color{techline}\leaders\hrule height \headrulewidth\hfill}}
\setlength{\parskip}{5pt}
\setlength{\parindent}{0pt}
\setlength{\emergencystretch}{3em}
\setlist{itemsep=3pt,topsep=4pt}
\captionsetup{font=small,labelfont={bf,color=techpurple},labelsep=colon}
\lstset{basicstyle=\ttfamily\small,breaklines=true,frame=single,rulecolor=\color{techline},backgroundcolor=\color{techlavender},showstringspaces=false,keywordstyle=\color{techpurple}\bfseries,commentstyle=\color{gray}}
\newtcolorbox{techbox}[1][]{colback=techlavender,colframe=techline,arc=3mm,boxrule=.6pt,left=7pt,right=7pt,top=6pt,bottom=6pt,#1}
\newtcolorbox{warningbox}{colback=techpink!7,colframe=techpink,arc=3mm,boxrule=.6pt,left=7pt,right=7pt,top=6pt,bottom=6pt}
\newcommand{\TechFigure}[4]{%
\begin{figure}[H]\centering
\includegraphics[width=\linewidth,height=#2,keepaspectratio]{#1}
\caption{#3}\label{#4}
\end{figure}}
\newcommand{\TechCover}[3]{%
\begin{titlepage}
\begin{tikzpicture}[remember picture,overlay]
\fill[techlavender] (current page.south west) rectangle (current page.north east);
\fill[techpurple] (current page.north west) rectangle ([yshift=-3.8cm]current page.north east);
\fill[techpurple2] ([xshift=-3.4cm]current page.south east) circle (4.4cm);
\draw[techpurple2,line width=2pt] ([xshift=1.5cm,yshift=-7.4cm]current page.north west) circle (2.2cm);
\draw[white,line width=1.5pt] ([xshift=1.5cm,yshift=-7.4cm]current page.north west) circle (1.45cm);
\end{tikzpicture}
\vspace*{0.65cm}\color{white}\raggedright
{\fontsize{24}{28}\selectfont\bfseries TECH-FLOW\par}
{\large Sistema de gestión de inventario y ventas\par}
\vspace{3.2cm}\color{techink}
{\fontsize{30}{35}\selectfont\bfseries #1\par}
\vspace{0.35cm}{\Large\color{techpurple}#2\par}
\vfill
\begin{techbox}[colback=white,colframe=techline]
\textbf{Universidad:} Universidad Mayor de San Simón\\
\textbf{Docente:} CORINA JUSTINA FLORES VILLARROEL\\
\textbf{Integrantes:}\\
Jorge Henry Quispe Escalera\\
Logan Betho Canelas Silvestre\\
Fausto Jafeth Vilches Mendieta\\
Eunice Melany Zurita Pérez\\
\textbf{Entrega:} 30 de noviembre de 2026\\
\textbf{Defensas:} 1, 2 y 8 de diciembre de 2026\\
\end{techbox}
\end{titlepage}}

\begin{document}
\TechCover{Manual de instalación}{Puesta en marcha de Tech-Flow en su equipo}{Manual de instalación}
\tableofcontents\clearpage

\section{Antes de comenzar}
Este manual explica cómo dejar Tech-Flow funcionando en su equipo, desde la obtención del código hasta la primera ejecución verificada. Está escrito para equipos con Windows 10/11 y Linux; en Linux se han verificado escritorios sencillos y gestores de ventanas en mosaico (bspwm), donde la ventana única encaja naturalmente.

\section{Requisitos técnicos}
\begin{table}[H]\centering\caption{Requisitos verificados}\begin{tabular}{ll}\toprule
Componente & Versión verificada\\\midrule
Java JDK & 17 (zulu-jdk-fx o equivalente con JavaFX)\\
Apache Maven & 3.x con acceso al repositorio central\\
JavaFX & 17.0.2 (vía Maven Central)\\
SO & Linux o Windows 10/11\\
Espacio libre & ~500 MB (incluido el repositorio local de Maven, \texttt{\textasciitilde/.m2})\\\bottomrule\end{tabular}\end{table}

Opcional: IntelliJ IDEA 2024+ o Eclipse 2023+ si prefiere un IDE con proyecto Maven importado. No se requiere base de datos ni servidor.

\section{1. Obtener el proyecto}
\begin{itemize}
\item Con Git: \texttt{git clone https://github.com/fausto-lang/Tech-Flow.git} y \texttt{cd Tech-Flow}.
\item Sin Git: descargue el ZIP del repositorio y extraiga el contenido en una carpeta sin espacios.
\end{itemize}

\section{2. Verificar Java y Maven}
\texttt{java -version} debe reportar 17 o superior (build 17). \texttt{mvn -v} debe mostrar Apache Maven. Si \texttt{mvn -v} falla, instale Maven (Linux: \texttt{sudo apt install maven} o el gestor de su distribución; Windows: oriente \texttt{PATH} a \texttt{\%MAVEN\_HOME\%\textbackslash bin}).

\section{3. Compilar y ejecutar pruebas}
\begin{itemize}
\item \texttt{mvn compile} — compila el proyecto completo.
\item \texttt{mvn test} — ejecuta la suite JUnit 5 (516/516 en verde).
\end{itemize}
La primera descarga de dependencias tarda más; las siguientes usan la caché de \texttt{.m2}.

\section{4. Ejecutar Tech-Flow}
\begin{itemize}
\item \texttt{mvn javafx:run} — arranca la ventana única en el \texttt{com.overbyte.Main} configurado en el \texttt{pom.xml}.
\item En Linux el proyecto y la documentación vienen adaptados a bspwm; la ventana se abre con transiciones suaves sin diálogo flotante.
\end{itemize}

\section{5. Verificar la instalación}
\begin{enumerate}
\item Se abre una única ventana con la pantalla de inicio de sesión.
\item Inicie sesión como administrador de demostración (\texttt{adm} / \texttt{user}).
\item Abra la caja (panel administrativo). El icono de caja cambia a estado abierto.
\item Vaya a Inventario y confirme que la tabla carga al menos una fila.
\item Cierre sesión y vuelva a entrar como cajero demo o almacén demo si fue provisto con su instancia.
\end{enumerate}

\section{Errores frecuentes}
\begin{itemize}
\item \textbf{``java: command not found''}: JDK no instalado o \texttt{JAVA\_HOME} no orientado al 17.
\item \textbf{``mvn: command not found''}: Maven no instalado o \texttt{PATH} incompleto.
\item \textbf{``module not found: javafx''}: ejecute \texttt{mvn javafx:run} y no \texttt{java} a secas; el plugin monta el \texttt{--module-path} correcto.
\item \textbf{``blanco/blanco'' y nada pasa al hacer clic en Entrar}: credenciales o rol incorrectos; el login exige los tres datos consistentes.
\item \textbf{``Acceso restringido''}: caja cerrada; un administrador debe abrirla antes de vender.
\end{itemize}

\section{Desinstalación}
Confirme que ningún proceso \texttt{java} de Tech-Flow está en ejecución, elimine la carpeta del proyecto y limpie \texttt{.m2/repository/org/openjfx} y las rutas de Maven si desea liberar la caché. Los datos en \texttt{data/} se borran junto con el proyecto; realice una copia si los necesita.

\end{document}
